\documentclass[11pt,a4paper]{article}

\usepackage[utf8]{inputenc}
\usepackage[T1]{fontenc}
\usepackage{lmodern}
\usepackage{amsmath,amssymb,amsthm}
\usepackage{booktabs}
\usepackage{hyperref}
\usepackage{geometry}
\usepackage{fancyhdr}
\usepackage{enumitem}
\usepackage{listings}
\usepackage{xcolor}
\usepackage{graphicx}
\usepackage{caption}

\geometry{margin=1in}

\hypersetup{
    colorlinks=true,
    linkcolor=blue,
    urlcolor=blue,
    citecolor=blue
}

\lstset{
    basicstyle=\ttfamily\small,
    breaklines=true,
    frame=single,
    backgroundcolor=\color{gray!10}
}

\title{\textbf{pqcoin: A Post-Quantum Cryptocurrency}}
\author{Version 0.1.0}
\date{}

\begin{document}

\maketitle

\begin{abstract}
pqcoin is a proof-of-work cryptocurrency designed to provide security against both classical and quantum computers. By utilizing NIST-standardized post-quantum cryptographic primitives---specifically ML-DSA-87 (FIPS 204) for digital signatures and SHA3-512 (FIPS 202) for hashing---pqcoin offers a quantum-resistant alternative to existing cryptocurrencies that rely on elliptic curve cryptography. The design follows Bitcoin's proven UTXO model while making deliberate simplifications: removing the scripting language in favor of two well-defined output types (P2PKH and M-of-N multisig), and modernizing the cryptographic foundation for the post-quantum era.
\end{abstract}

\tableofcontents
\newpage

\section{Introduction}

\subsection{Motivation}

Current cryptocurrencies like Bitcoin rely on the Elliptic Curve Digital Signature Algorithm (ECDSA) for transaction authorization. While secure against classical computers, ECDSA is vulnerable to Shor's algorithm running on a sufficiently powerful quantum computer. As quantum computing advances, the cryptocurrency ecosystem faces an existential threat to its security model.

pqcoin addresses this threat by building on cryptographic primitives that are believed to be secure against both classical and quantum attacks, as standardized by NIST in their Post-Quantum Cryptography project.

\subsection{Design Goals}

\begin{enumerate}
    \item \textbf{Quantum Resistance}: Use NIST-standardized post-quantum cryptographic algorithms
    \item \textbf{Simplicity}: Remove unnecessary complexity while maintaining Bitcoin's proven security model
    \item \textbf{Correctness}: Prioritize security and correctness over feature richness
    \item \textbf{Compatibility}: Follow familiar patterns from Bitcoin where appropriate
\end{enumerate}

\section{Cryptographic Primitives}

\subsection{Hash Function: SHA3-512 (FIPS 202)}

All hashing operations in pqcoin use SHA3-512, providing:

\begin{itemize}
    \item \textbf{Output size}: 512 bits (64 bytes)
    \item \textbf{Security level}: 256-bit classical security
    \item \textbf{Quantum security}: $\sim$256-bit security against Grover's algorithm (effectively 128-bit due to birthday attacks)
\end{itemize}

SHA3-512 is used for:
\begin{itemize}
    \item Transaction ID computation
    \item Block header hashing (proof-of-work)
    \item Address derivation (hash of public key)
    \item Merkle tree construction
\end{itemize}

\subsection{Digital Signatures: ML-DSA-87 (FIPS 204)}

Transaction authorization uses ML-DSA-87 (formerly Dilithium5), a lattice-based digital signature scheme. Table~\ref{tab:mldsa} summarizes the key sizes.

\begin{table}[h]
\centering
\caption{ML-DSA-87 Key and Signature Sizes}
\label{tab:mldsa}
\begin{tabular}{@{}lr@{}}
\toprule
Component & Size (bytes) \\
\midrule
Public Key & 2,592 \\
Secret Key & 4,896 \\
Signature & 4,627 \\
\bottomrule
\end{tabular}
\end{table}

ML-DSA-87 provides NIST Level 5 security ($\sim$256-bit classical, $\sim$128-bit quantum security). This is the highest security level in the ML-DSA family, providing maximum protection against future advances in quantum computing.

\subsection{Signature Construction}

Signatures are computed over the SHA3-512 hash of the transaction's signing data:

\begin{lstlisting}
message = SHA3-512(signing_data(tx, input_index))
signature = ML-DSA-87.Sign(secret_key, message)
\end{lstlisting}

This approach:
\begin{enumerate}
    \item Provides a fixed-size input regardless of transaction size
    \item Matches Bitcoin's practice of signing transaction hashes
    \item Is compatible with ML-DSA-87's design for signing arbitrary messages
\end{enumerate}

\section{Transaction Model}

\subsection{UTXO Model}

pqcoin uses the Unspent Transaction Output (UTXO) model, identical in concept to Bitcoin:

\begin{itemize}
    \item Transactions consume existing outputs (inputs) and create new outputs
    \item Each output can only be spent once
    \item The sum of input values must equal or exceed output values
    \item The difference (inputs $-$ outputs) is the transaction fee
\end{itemize}

\subsection{Output Types}

Unlike Bitcoin's scripting system, pqcoin supports exactly two output types:

\subsubsection{Pay-to-Public-Key-Hash (P2PKH)}

The standard single-signature output. Funds are locked to the SHA3-512 hash of a public key. To spend:
\begin{enumerate}
    \item Provide the public key that hashes to the address
    \item Provide a valid signature over the transaction
\end{enumerate}

\subsubsection{M-of-N Multisig}

Multi-signature outputs requiring $M$ signatures from $N$ possible signers. Limited to a maximum of 15 keys due to post-quantum signature sizes.

\subsection{Transaction Structure}

\begin{lstlisting}
Transaction {
    version: u32,           // Currently 1
    inputs: Vec<TxInput>,   // References to outputs being spent
    outputs: Vec<TxOutput>, // New outputs being created
}

TxInput {
    outpoint: OutPoint,     // (txid, index) reference
    witness: Witness,       // Proof of authorization
}

TxOutput {
    amount: u64,            // Value in quanta
    condition: LockingCondition,  // P2PKH or Multisig
}
\end{lstlisting}

\subsection{Dust Limit}

To prevent UTXO set bloat, outputs must exceed a dynamic dust limit based on the cost to spend them:

\begin{equation}
\text{dust\_limit} = \frac{\text{estimated\_spend\_size} \times \text{DUST\_FEE\_RATE}}{1000}
\end{equation}

Where DUST\_FEE\_RATE = 100 quanta per KB. This ensures outputs are economically rational to spend, accounting for the larger post-quantum signatures.

\section{Block Structure}

\subsection{Block Header}

The block header is 176 bytes, structured as shown in Table~\ref{tab:header}.

\begin{table}[h]
\centering
\caption{Block Header Structure}
\label{tab:header}
\begin{tabular}{@{}llp{7cm}@{}}
\toprule
Field & Size & Description \\
\midrule
version & 4 bytes & Block format version \\
prev\_hash & 64 bytes & SHA3-512 hash of previous header \\
merkle\_root & 64 bytes & Root of transaction merkle tree \\
timestamp & 8 bytes & Unix timestamp \\
difficulty\_bits & 4 bytes & Compact difficulty target \\
nonce & 32 bytes & Proof-of-work nonce \\
\bottomrule
\end{tabular}
\end{table}

\subsection{256-bit Nonce}

pqcoin uses a 32-byte (256-bit) nonce, providing $2^{256}$ possible values. This design:
\begin{itemize}
    \item Eliminates nonce exhaustion concerns
    \item Removes the need for Bitcoin-style extraNonce mechanisms
    \item Simplifies mining implementation
\end{itemize}

\subsection{Merkle Root}

The merkle root commits to all transactions using a binary hash tree:
\begin{itemize}
    \item Leaves are transaction IDs (SHA3-512 hashes)
    \item Internal nodes are SHA3-512 hashes of concatenated children
    \item Odd levels duplicate the last node
\end{itemize}

\subsection{Block Limits}

\begin{table}[h]
\centering
\caption{Block Limits}
\label{tab:blocklimits}
\begin{tabular}{@{}lr@{}}
\toprule
Parameter & Value \\
\midrule
Max transactions & 1,000 (excluding coinbase) \\
Max block size & 16 MB \\
Max inputs per tx & 10,000 \\
Max outputs per tx & 10,000 \\
Max multisig keys & 15 \\
\bottomrule
\end{tabular}
\end{table}

The 16 MB block size (vs Bitcoin's effective $\sim$4 MB) accommodates the larger post-quantum signatures.

\section{Consensus Mechanism}

\subsection{Proof of Work}

Valid blocks must have a header hash less than or equal to the target:

\begin{equation}
\text{SHA3-512}(\text{block\_header}) \leq \text{target}
\end{equation}

The target is derived from difficulty\_bits using the formula:
\begin{equation}
\text{target} = \text{coefficient} \times 2^{8 \times (\text{exponent} - 3)}
\end{equation}

Where $\text{difficulty\_bits} = (\text{exponent} \ll 24) \mid \text{coefficient}$.

\subsection{Difficulty Adjustment}

Difficulty adjusts every 2,016 blocks to maintain a 10-minute average block time:

\begin{enumerate}
    \item Calculate actual time for the previous 2,016 blocks
    \item Clamp adjustment to 4$\times$ in either direction
    \item Scale the target: $\text{new\_target} = \text{old\_target} \times \frac{\text{actual\_time}}{\text{target\_time}}$
\end{enumerate}

A minimum difficulty floor (MIN\_DIFFICULTY\_BITS = 0x41ffffff) prevents difficulty from dropping below a safe threshold.

\subsection{Chain Selection}

The valid chain with the most cumulative proof-of-work is the canonical chain. Reorganizations are limited to 100 blocks to prevent long-range attacks.

\subsection{Timestamp Rules}

Block timestamps must:
\begin{enumerate}
    \item Be strictly greater than the Median-Time-Past (median of previous 11 blocks)
    \item Not exceed current time + 2 hours
\end{enumerate}

This prevents timestamp manipulation attacks near difficulty adjustment boundaries.

\section{Monetary Policy}

\subsection{Block Reward}

\begin{table}[h]
\centering
\caption{Monetary Parameters}
\label{tab:monetary}
\begin{tabular}{@{}lr@{}}
\toprule
Parameter & Value \\
\midrule
Initial reward & 50,000,000 quanta (50 coins) \\
Halving interval & 210,000 blocks \\
Smallest unit & 1 quantum \\
\bottomrule
\end{tabular}
\end{table}

\subsection{Supply Schedule}

Following Bitcoin's model:
\begin{itemize}
    \item Total supply asymptotically approaches 21 million coins
    \item Reward halves every $\sim$4 years (at 10-minute blocks)
    \item After 64 halvings, reward becomes 0
\end{itemize}

\subsection{Coinbase Maturity}

Coinbase outputs cannot be spent until 100 blocks have been mined on top of them. This prevents issues from chain reorganizations affecting miner rewards.

\section{Network Protocol}

\subsection{Message Format}

Messages use a length-prefixed framing format:

\begin{lstlisting}
[Magic: 4 bytes "PQCN"]
[Command: 12 bytes]
[Length: 4 bytes]
[Checksum: 4 bytes]
[Payload: variable]
\end{lstlisting}

The checksum is the first 4 bytes of SHA3-512(payload).

\subsection{Message Types}

\begin{table}[h]
\centering
\caption{Network Message Types}
\label{tab:messages}
\begin{tabular}{@{}ll@{}}
\toprule
Message & Purpose \\
\midrule
version/verack & Handshake and capability exchange \\
ping/pong & Keepalive \\
inv/getdata & Inventory announcement and request \\
getheaders/headers & Header synchronization \\
block/tx & Block and transaction relay \\
addr/addrv2 & Peer discovery \\
\bottomrule
\end{tabular}
\end{table}

\subsection{Peer Management}

\begin{table}[h]
\centering
\caption{Network Parameters}
\label{tab:network}
\begin{tabular}{@{}lr@{}}
\toprule
Parameter & Value \\
\midrule
Default port & 8333 \\
Max connections & 125 (8 outbound) \\
Max per IP & 3 \\
Max per /16 subnet & 2 \\
Ban duration & 24 hours \\
Ping interval & 60 seconds \\
\bottomrule
\end{tabular}
\end{table}

Subnet limits help prevent eclipse attacks.

\section{Wallet}

\subsection{Key Storage}

Wallet files use encrypted JSON format:
\begin{itemize}
    \item Secret keys encrypted with AES-256-GCM
    \item Key derivation via Argon2 from user password
    \item Memory is zeroized after use
\end{itemize}

\subsection{Address Format}

Addresses are the full 64-byte SHA3-512 hash of the public key, displayed as 128 hexadecimal characters.

\section{Security Considerations}

\subsection{Quantum Resistance}

pqcoin's security against quantum computers rests on:
\begin{enumerate}
    \item \textbf{ML-DSA-87}: Based on the Module-LWE problem, believed to be hard for quantum computers
    \item \textbf{SHA3-512}: Grover's algorithm provides only quadratic speedup, leaving 256-bit security
\end{enumerate}

\subsection{Classical Security}

Against classical computers:
\begin{itemize}
    \item 256-bit hash security (SHA3-512)
    \item 256-bit signature security (ML-DSA-87 at NIST Level 5)
\end{itemize}

\subsection{Known Limitations}

\begin{enumerate}
    \item \textbf{Signature size}: ML-DSA-87 signatures (4,627 bytes) are $\sim$72$\times$ larger than ECDSA (64 bytes), affecting transaction throughput
    \item \textbf{No scripting}: Complex smart contracts are not supported
    \item \textbf{No chained transactions}: Mempool does not support spending unconfirmed outputs
\end{enumerate}

\section{Comparison with Bitcoin}

\begin{table}[h]
\centering
\caption{pqcoin vs Bitcoin}
\label{tab:comparison}
\begin{tabular}{@{}lll@{}}
\toprule
Feature & pqcoin & Bitcoin \\
\midrule
Hash function & SHA3-512 (64 bytes) & SHA-256 (32 bytes) \\
Signature scheme & ML-DSA-87 & ECDSA (secp256k1) \\
Signature size & 4,627 bytes & 64 bytes \\
Public key size & 2,592 bytes & 33 bytes \\
Quantum resistant & Yes & No \\
Scripting & No (P2PKH + Multisig) & Yes (Script) \\
Nonce size & 256 bits & 32 bits \\
Block size limit & 16 MB & $\sim$4 MB \\
\bottomrule
\end{tabular}
\end{table}

\section{Future Work}

\begin{enumerate}
    \item \textbf{Light clients}: Bloom filtering for SPV wallets
    \item \textbf{Checkpoints}: Hardcoded block hashes for faster initial sync
    \item \textbf{Assume-valid}: Skip signature verification for deeply-buried blocks
    \item \textbf{Hardware wallet support}: Integration with post-quantum capable hardware
\end{enumerate}

\section{Conclusion}

pqcoin provides a quantum-resistant cryptocurrency built on NIST-standardized cryptographic primitives. By simplifying Bitcoin's design while upgrading its cryptographic foundation, pqcoin offers a secure platform for value transfer in the post-quantum era.

The deliberate constraints---no scripting, limited output types, no chained transactions---prioritize security and correctness over feature richness, making the codebase more auditable and the security model more analyzable.

\section*{References}

\begin{enumerate}
    \item NIST FIPS 202: SHA-3 Standard: Permutation-Based Hash and Extendable-Output Functions
    \item NIST FIPS 204: Module-Lattice-Based Digital Signature Standard
    \item Nakamoto, S. (2008). Bitcoin: A Peer-to-Peer Electronic Cash System
    \item NIST Post-Quantum Cryptography Standardization Process
\end{enumerate}

\appendix

\section{Protocol Constants}

\begin{lstlisting}
// Protocol
MAX_BLOCK_TXS = 1,000
MAX_TX_INPUTS = 10,000
MAX_TX_OUTPUTS = 10,000
MAX_MULTISIG_KEYS = 15
MAX_BLOCK_SIZE = 16 MB
COINBASE_MATURITY = 100 blocks
MAX_REORG_DEPTH = 100 blocks
MAX_FUTURE_BLOCK_TIME = 2 hours

// Mining
INITIAL_REWARD = 50,000,000 quanta
DIFFICULTY_INTERVAL = 2,016 blocks
TARGET_BLOCK_TIME = 600 seconds
HALVING_INTERVAL = 210,000 blocks
MIN_DIFFICULTY_BITS = 0x41ffffff

// Mempool
MIN_RELAY_FEE = 1,000 quanta/byte
DUST_FEE_RATE = 100 quanta/KB
TX_EXPIRY_BLOCKS = 432 blocks (~72 hours)
MAX_MEMPOOL_SIZE = 5,000 transactions

// Network
DEFAULT_PORT = 8333
MAX_PEERS = 125
MAX_OUTBOUND = 8
BAN_DURATION = 24 hours
PING_INTERVAL = 60 seconds
\end{lstlisting}

\section{Cryptographic Key Sizes}

\begin{table}[h]
\centering
\caption{Cryptographic Component Sizes}
\label{tab:keysizes}
\begin{tabular}{@{}llr@{}}
\toprule
Algorithm & Component & Size (bytes) \\
\midrule
SHA3-512 & Hash output & 64 \\
ML-DSA-87 & Public key & 2,592 \\
ML-DSA-87 & Secret key & 4,896 \\
ML-DSA-87 & Signature & 4,627 \\
AES-256-GCM & Key & 32 \\
AES-256-GCM & Nonce & 12 \\
Argon2 & Salt & 16 \\
\bottomrule
\end{tabular}
\end{table}

\end{document}
